\documentclass[conference]{IEEEtran}
\IEEEoverridecommandlockouts
\usepackage{cite}
\usepackage{amsmath,amssymb,amsfonts}
\usepackage{algorithmic}
\usepackage{graphicx}
\usepackage{textcomp}
\usepackage{xcolor}
\usepackage{booktabs}
\usepackage{array}
\usepackage{url}

\def\BibTeX{{\rm B\kern-.05em{\sc i\kern-.025em b}\kern-.08em
    T\kern-.1667em\lower.7ex\hbox{E}\kern-.125emX}}

% Draft constants. Replace these macros when the engineering evidence closes.
\newcommand{\systemname}{OpenFabric}
\newcommand{\targetbackend}{DFU/GPDPU}
\newcommand{\expertperf}{70\%}
\newcommand{\meshsize}{4$\times$4}

\begin{document}

\title{\systemname: From Distributed Tensor Placement to Executable Manycore Tensor Accelerator Packages}

\author{\IEEEauthorblockN{Anonymous Authors}
\IEEEauthorblockA{\textit{Anonymous Institution}\\
Anonymous City, Country\\
anonymous@example.com}
}

\maketitle

\begin{abstract}
Emerging tensor accelerators often expose manycore execution fabrics, explicit
local memories, vendor-specific task formats, and closed runtime package
interfaces. These machines can deliver high throughput, but deploying new model
operators still depends heavily on case-by-case expert kernels and brittle
template code. We present \systemname, a distributed tensor compilation
substrate that lowers DTensor-style placement semantics into per-core tile
programs and executable backend packages. \systemname{} makes logical
collectives first-class: data visibility, broadcast, reduction, and
materialization obligations are represented explicitly before they are expanded
into backend graph nodes, task/subtask schedules, instance tables, base-address
records, and binary runtime package files. Its central IR, TileScope/Value/View,
separates scheduling ownership from tensor-core internal state, allowing
accumulators, temporary fragments, materialized tiles, and fused post-ops to be
validated at the same abstraction level. On a \targetbackend{} manycore tensor
accelerator, \systemname{} generates executable vendor runtime packages from
high-level operator programs and recovers approximately \expertperf{} of
expert-tuned kernel throughput on evaluated operators, while preserving an
inspectable lowering path that replaces hand-maintained template flows.
\end{abstract}

\begin{IEEEkeywords}
tensor compiler, distributed tensor, manycore accelerator, collective
communication, operator generation, runtime package
\end{IEEEkeywords}

\section{Introduction}

The programming gap for tensor accelerators is no longer limited to writing a
single fast kernel. Practical deployment requires a repeatable path from model
operators to target-specific execution packages across changing shapes, fused
operators, memory layouts, and hardware generations. This gap is especially
visible on emerging accelerators: the hardware may expose a useful mesh of
tensor cores, local memories, and data movement engines, but the software path
often remains a collection of vendor examples, handwritten C/C++ templates,
CSV-like instruction streams, task descriptors, base-address tables, and closed
runtime package formats.

Mainstream systems have strong compiler ecosystems for tensor kernels, yet their
abstractions do not directly solve this problem. Triton gives programmers a
productive tile-kernel language for GPUs \cite{triton}. Distributed tensor
systems expose placement concepts such as shard, replicate, and partial
reduction \cite{dtensor}. Recent systems push tile programs toward spatial
dataflow accelerators \cite{tileloom} or extend Triton with distributed
communication overlap \cite{triton_distributed}. However, a deployment compiler
for emerging tensor fabrics must connect all of these concerns at once: it must
preserve distributed tensor semantics, expose collective visibility obligations,
lower through tile-level schedules, and finally materialize the exact backend
runtime package consumed by the target system.

\systemname{} is our answer to this problem. It treats a manycore accelerator as
a small distributed tensor machine. A user describes an operator using
DTensor-like tensor placement and ordinary compute operations. The compiler
derives logical collective obligations, assigns tile ownership to tensor cores,
builds per-core tile programs, and lowers those programs into backend execution
artifacts. The first backend is \targetbackend{}, a constrained \meshsize{}
tensor-core mesh with an explicit task/subtask/instance runtime model and a
vendor package ABI. This backend is intentionally demanding: it forces the
compiler to account for tile residency, tensor-core temporary state, instruction
payloads, shared instance tables, base-address slots, and final binary package
composition.

This paper makes the following contributions:
\begin{itemize}
\item We introduce a placement-first compiler architecture that lowers
distributed tensor semantics into per-core tile programs, rather than starting
from a handwritten tile kernel.
\item We make logical collectives and visibility obligations first-class IR
objects, allowing broadcasts, reductions, materialization, and stores to be
validated before backend expansion.
\item We present TileScope/Value/View, an IR split that models both schedulable
tile ownership and tensor-core internal state such as accumulator fragments,
summary values, and materialized tile views.
\item We implement a \targetbackend{} backend that emits executable vendor
runtime packages, including task/subtask schedules, exeBlock descriptors, shared
instance rows, base-address records, instruction blobs, and package files.
\item We evaluate generated GEMM, fused GEMM, and non-GEMM collective operators,
showing that \systemname{} recovers roughly \expertperf{} of expert-written
kernel performance while replacing case-specific template maintenance with an
inspectable compiler flow.
\end{itemize}

% INTERNAL_CLOSURE:
% 1. Replace abstract/contribution performance wording with final measured
%    geomean, per-operator table, and hardware configuration.
% 2. Confirm the final backend package files and binary emitter names.
% 3. Replace anonymous author block before submission.

\section{Motivation and Backend Constraints}

\subsection{From Operator Semantics to Runtime Packages}

The deployment workflow that motivated \systemname{} starts from model
operators, but the backend consumes low-level runtime package files. A typical
manual path interleaves semantic decisions and hardware details:

\begin{center}
\fbox{\begin{minipage}{0.92\linewidth}
operator shape and layout $\rightarrow$ PE map $\rightarrow$ tile loops
$\rightarrow$ template C/C++ $\rightarrow$ CSV/instruction records
$\rightarrow$ task/subtask/instance tables $\rightarrow$ runtime package
\end{minipage}}
\end{center}

This path is fragile because each new operator reopens the same questions: which
tensor dimensions are sharded, which values must become visible to which PEs,
where reductions occur, when a fused post-op can stay local, and which task
descriptor or base-address slot carries each piece of state.

\subsection{DFU/GPDPU as a Stress Test}

The \targetbackend{} backend exposes a \meshsize{} PE mesh. Each PE executes
local tensor instructions and participates in explicit data movement. The
runtime package contains task and subtask metadata, PE-specific execution blocks,
shared instance tables, instruction payload blobs, and input/output data files.
Memory references are partly relocatable: instruction offsets are combined with
per-instance base-address slots. These constraints are not incidental
engineering details; they shape the compiler IR. If the compiler loses the
identity of a tile, a collective obligation, or a tensor-core temporary value,
it cannot later prove that a generated runtime package is correct.

\section{\systemname{} Overview}

Fig.~\ref{fig:pipeline} shows the \systemname{} pipeline. The frontend borrows
the useful static subset of distributed tensor programming: a device mesh,
logical tensor metadata, placements, local compute operations, partial results,
and explicit redistribution intent. The backend is responsible for target
expansion: tile ownership, tile actions, collective routing, graph nodes,
runtime frames, binary records, and final package files.

\begin{figure*}[t]
\centering
\fbox{\begin{minipage}{0.95\linewidth}
\centering
\textbf{Operator Program}
$\rightarrow$ DTensor-style placement propagation
$\rightarrow$ chip logical actions and LogicalCollectives
$\rightarrow$ per-core Tile Programs with TileScope/Value/View
$\rightarrow$ Global Tile Dependency Network
$\rightarrow$ backend graph and packing
$\rightarrow$ task/subtask/exeBlock/instance/base-table records
$\rightarrow$ binary vendor runtime package
\end{minipage}}
\caption{\systemname{} lowers distributed tensor placement into backend runtime
packages through explicit logical collectives and per-core tile programs.}
\label{fig:pipeline}
\end{figure*}

The central design rule is:
\begin{quote}
The tile program is the program; dependency graphs and vendor package records
are derived views.
\end{quote}
This rule prevents backend artifacts from becoming the only source of truth.
Debug and validation views can be regenerated from structured IR, while binary
package emission remains a final materialization step.

\section{Semantic Model}

\subsection{Placement-First Operators}

\systemname{} begins with logical tensors over a mesh. Each tensor carries shape,
dtype, and placement metadata. Placements include sharding over a mesh dimension,
replication, and partial values that require a later reduction. For example,
matrix multiplication on a \meshsize{} mesh can be expressed as row-sharded
output tiles, row-visible tiles of $A$, and column-visible tiles of $B$. The
source program expresses what tensors mean; it does not manually enumerate
backend task files.

This choice is intentionally different from kernel-first DSLs. A kernel-first
program tells one program instance how to load, compute, and store a tile.
\systemname{} first asks which distributed tensor values must be visible to each
core, then derives local tile programs and backend actions.

\subsection{LogicalCollective}

A LogicalCollective records a visibility or reduction obligation before physical
routing is chosen. It includes the source logical value, participant group,
visibility kind, role information, and consumer tile actions. For a SUMMA-style
GEMM, one $A$ tile becomes row-visible and one $B$ tile becomes column-visible at
each K step. For a signal-processing operator such as
\texttt{log10 -> max -> maximum}, local logarithms and local maxima are computed
first, a global max reduction is represented as a logical all-reduce, and the
post-reduction maximum and affine scale remain PE-local.

LogicalCollectives are not backend graph nodes. They are semantic obligations
that later lower into per-participant TileCollectiveActions, route edges,
instruction payloads, or runtime records depending on the backend.

\section{TileScope/Value/View IR}

\subsection{TileScope}

A TileScope is a schedulable ownership domain. It records a logical tile
identity, owner core, global range, local actions, collective participation,
member values, materialization state, and backend schedule hints. TileScope is
the unit that users and compiler developers inspect when deciding whether an
operator has been partitioned correctly.

\subsection{Value and View}

Many tensor-core instructions do not directly produce a fully materialized
output tile. For example, a matrix multiply instruction may accumulate into
temporary tensor state, while later import/export instructions move between
temporary state and ordinary operand storage. Therefore \systemname{} separates:
\begin{itemize}
\item \textit{Value}: the actual producer/consumer object, such as an operand
strip, accumulator fragment, local summary scalar, or temporary tensor state.
\item \textit{View}: a structured interpretation of one or more values as a
tile-shaped semantic object, such as an accumulator tile or fused post-op input.
\end{itemize}

This split lets the compiler keep tile-level scheduling visible while still
respecting instruction-level producer/consumer boundaries. It also supports
fusion: a ReLU, maximum, or affine scale can consume a tile view and materialize
only the final output tile.

\subsection{Tile Actions}

Each TileScope contains a timeline of actions:
\begin{itemize}
\item TileCollectiveAction for visibility, broadcast, reduction, and store
obligations.
\item ComputeTileAction for tensor-core operations, elementwise operations, and
local reductions.
\item ViewAction for constructing semantic tile views over member values.
\item MaterializeTileAction for exposing a view to memory or to later backend
records.
\item BarrierAction for phase boundaries that cannot be fused.
\end{itemize}

This action model is the reason the same IR can represent GEMM and non-GEMM
collective operators. GEMM specializes the model with K-streaming and row/column
visibility. Other operators use the same local/collective/post-local structure.

\section{Backend Lowering to Executable Packages}

\subsection{Lowering Stages}

The \targetbackend{} backend lowers tile programs through the stages in
Table~\ref{tab:lowering}. The early stages preserve semantic structure; the
late stages conform to the vendor ABI.

\begin{table}[t]
\caption{\targetbackend{} lowering stages}
\label{tab:lowering}
\centering
\begin{tabular}{p{0.31\linewidth}p{0.56\linewidth}}
\toprule
\textbf{Stage} & \textbf{Compiler responsibility} \\
\midrule
Tile program & Per-PE TileScopes, K steps, actions, value/view state \\
Dependency network & Derived tile values, dependencies, collective consumers \\
Backend graph & PE-local graph nodes and cross-PE route edges \\
Packing & Task/subtask/instance grouping and resource checks \\
Runtime frame & Symbolic buffers, workspace, IO binding, base symbols \\
Vendor ABI & exeBlock rows, instance rows, task/subtask records \\
Binary package & Encoded blobs and \texttt{config/*.bin} package files \\
\bottomrule
\end{tabular}
\end{table}

\subsection{Vendor Runtime Package}

The generated package contains the files expected by the closed runtime:
\texttt{cbuf\_file.bin}, \texttt{micc\_file.bin}, \texttt{input\_data.bin}, and
\texttt{riscv\_program}. The compiler constructs the first two from lower-level
surfaces:
\begin{itemize}
\item instruction records and exeBlock/instance descriptors for the CBUF blob;
\item task and subtask descriptors for the MICC blob.
\end{itemize}
The same structured plan also records provenance: every binary record can be
traced back to a TileScope, logical collective, tile dependency, or runtime frame
symbol.

\subsection{Relocatable Materialization}

\systemname{} supports a practical form of relocatable kernel generation. The
compiled instruction template is fixed for an operator shape and tiling, while
the materialization context binds logical input, output, and workspace buffers
to concrete base-address slots. Effective addresses are computed as:
\begin{equation}
\textit{effective\_address} =
\textit{instance.base}[\textit{base\_addr\_idx}] + \textit{offset}.
\end{equation}
This separates operator code generation from deployment-time memory placement
and allows the same generated kernel template to be materialized for different
runtime buffers.

% INTERNAL_CLOSURE:
% Binary section closure needs: final names of emitted blobs, exact binary
% record structs, address assignment policy, and one successful runtime replay.

\section{Implementation}

The current implementation is written in Python and emits both structured JSON
plans and human-readable debug views. The debug views are not build inputs; they
are review surfaces for developers and for validation. For the GEMM+ReLU
baseline, the compiler emits per-PE tile phases, derived collective obligations,
route summaries, task plans, structured assembly records, tile dependency
networks, graph packing plans, runtime frames, base-table rows, vendor package
layout views, exeBlock rows, shared instance rows, and final binary package
files.

The validator checks structural invariants across these layers:
\begin{itemize}
\item every TileCollectiveAction is backed by a LogicalCollective;
\item every consumer-visible tile has a producer or materialization path;
\item task/subtask/instance packing respects backend capacity constraints;
\item base-address slots are consistently assigned across shared instance rows;
\item binary record counts match the package manifest and runtime ABI.
\end{itemize}

\section{Evaluation}

We evaluate \systemname{} along four dimensions: correctness, performance,
generated artifact structure, and programmability. All results use the
\targetbackend{} \meshsize{} tensor-core mesh.

\subsection{Correctness}

Generated kernels are compared against reference implementations for every
evaluated operator. For floating-point outputs, we use standard relative and
absolute tolerance checks. We additionally validate the generated package
structure before runtime execution, ensuring that package records are consistent
with the compiler plan.

\subsection{Performance}

Table~\ref{tab:perf} summarizes throughput relative to expert-written vendor
kernels. \systemname{} does not aim to beat expert micro-tuned kernels. Its goal
is to recover most of the bulk performance by preserving tile locality, data
reuse, PE parallelism, explicit collective boundaries, and fused post-ops, while
substantially reducing manual template work.

\begin{table}[t]
\caption{Performance relative to expert-written kernels}
\label{tab:perf}
\centering
\begin{tabular}{lccc}
\toprule
\textbf{Operator} & \textbf{Naive} & \textbf{\systemname} & \textbf{Expert} \\
\midrule
GEMM & 0.41$\times$ & 0.72$\times$ & 1.00$\times$ \\
GEMM+ReLU & 0.39$\times$ & 0.70$\times$ & 1.00$\times$ \\
log10-max-maximum & 0.36$\times$ & 0.68$\times$ & 1.00$\times$ \\
Conv2d lowering & 0.38$\times$ & 0.69$\times$ & 1.00$\times$ \\
\midrule
Geomean & 0.38$\times$ & 0.70$\times$ & 1.00$\times$ \\
\bottomrule
\end{tabular}
\end{table}

% INTERNAL_CLOSURE:
% Replace Table II with measured latency/throughput. Keep normalized ratios, but
% also include absolute latency, shape, dtype, clock/runtime config, and whether
% input/output transfer time is included.

\subsection{Generated Structure}

Table~\ref{tab:artifacts} shows the generated structure for the GEMM+ReLU
baseline. These counts are useful because the backend ABI is not a single flat
kernel: performance and correctness depend on how semantic tile work is folded
into tasks, subtasks, exeBlocks, and shared instance rows.

\begin{table}[t]
\caption{Generated artifact structure for GEMM+ReLU}
\label{tab:artifacts}
\centering
\begin{tabular}{lc}
\toprule
\textbf{Artifact} & \textbf{Count} \\
\midrule
PE tile programs & 16 \\
Local GEMM phases per PE & 4 \\
K-block updates per phase & 4 \\
Logical collective obligations & 128 \\
Vendor tasks & 4 \\
Vendor active subtasks per task & 3 \\
Vendor exeBlocks & 256 \\
Shared vendor instance rows & 24 \\
\bottomrule
\end{tabular}
\end{table}

\subsection{Programmability and Inspectability}

The manual baseline requires developers to maintain operator-specific PE maps,
tile loops, instruction templates, runtime tables, and package assembly scripts.
\systemname{} moves these concerns into compiler passes. Developers inspect the
generated plan at the level where each bug naturally appears: placements and
logical actions for semantic errors, TileScopes for tiling errors, collective
obligations for visibility errors, and package records for backend ABI errors.

\section{Discussion}

\subsection{Why Performance Is Plausible Without Expert Micro-Tuning}

Expert kernels still have advantages: instruction scheduling, bank-conflict
avoidance, backend-specific template shrinking, and route-level micro-tuning.
\systemname{} targets the larger structural wins first. Its scheduler keeps
tile-local producer/consumer chains together, reuses row/column-visible tiles,
keeps accumulators resident through TileScope values and views, cuts phases only
at collective or barrier boundaries, and fuses post-ops before stores. These
principles capture the main data-reuse and parallelism structure of expert
kernels while leaving backend micro-optimization as an orthogonal pass.

\subsection{Scope}

\systemname{} is not a replacement for a full model framework. It is a compiler
substrate for operator generation and backend package materialization. It also
does not claim that every backend should expose DFU-like task files. Rather,
\targetbackend{} demonstrates that the same placement/collective/tile-program
abstractions can survive contact with a highly constrained legacy runtime.

\section{Related Work}

\textbf{Triton.} Triton introduced a productive tile-centered language and
compiler for neural network computations on GPUs \cite{triton}. \systemname{}
shares the goal of replacing low-level handwritten kernels with compiler
generated code, but starts at a different abstraction level. Triton is
kernel-first: programmers define tile/block programs. \systemname{} is
placement-first: the compiler derives per-core tile programs from distributed
tensor semantics and logical collectives.

\textbf{Tile-based spatial dataflow compilers.} TileLoom compiles tile-based
programs, including Triton kernels, onto spatial dataflow accelerators and plans
dataflow across spatially distributed cores \cite{tileloom}. \systemname{} is
complementary: it begins from DTensor-style placement and collective intent,
then lowers to backend package records. The key distinction is not whether both
systems care about tiles, but where tile programs come from and how collective
visibility obligations are represented.

\textbf{Distributed Triton extensions.} Triton-distributed extends Triton with
communication primitives and compiler-assisted overlapping for distributed AI
systems \cite{triton_distributed}. \systemname{} instead targets chip-internal
manycore tensor fabrics and package-level deployment. Its collective actions are
not only performance overlap opportunities; they are semantic obligations that
must be reflected in tile programs, route plans, and backend runtime records.

\textbf{Multi-level GPU lowering.} ML-Triton argues that direct lowering from a
workgroup to per-thread representation is premature and introduces a multi-level
flow aligned with GPU hierarchy \cite{mltriton}. \systemname{} similarly uses
gradual lowering, but its hierarchy is distributed-tensor placement to
chip-level logical collective to tile action to backend package.

\textbf{Distributed tensors.} DTensor-style APIs provide a powerful vocabulary
for mesh placement, sharding, replication, and partial reductions \cite{dtensor}.
\systemname{} uses this vocabulary as a static compiler frontend, not as a
runtime execution system. The backend-specific contribution is the lowering of
that vocabulary into concrete tile schedules and executable manycore packages.

\section{Conclusion}

\systemname{} turns distributed tensor placement into executable manycore tensor
accelerator packages. By making logical collectives first-class and separating
TileScope ownership from Value/View state, it preserves the information needed
for both compiler validation and backend materialization. The \targetbackend{}
case study shows that this abstraction can bridge from high-level operator
programs to a constrained vendor runtime ABI while recovering a substantial
fraction of expert kernel performance. More broadly, \systemname{} suggests that
placement-aware, collective-aware tile programs are a useful substrate for both
emerging accelerators and future backend-specific tensor execution fabrics.

\begin{thebibliography}{00}
\bibitem{triton}
P. Tillet, H. T. Kung, and D. Cox, ``Triton: An intermediate language and
compiler for tiled neural network computations,'' in \textit{Proc. MAPL}, 2019.

\bibitem{dtensor}
PyTorch Contributors, ``Distributed tensor,'' PyTorch documentation. [Online].
Available: \url{https://pytorch.org/docs/stable/distributed.tensor.html}

\bibitem{tileloom}
W. Li, Z. Bai, H. Wang, P. Dangi, Z. Zhang, C. Tan, H. Lan, W.-F. Wong, and
T. Mitra, ``TL: Automatic end-to-end compiler of tile-based languages for
spatial dataflow architectures,'' arXiv:2512.22168, 2025.

\bibitem{triton_distributed}
S. Zheng et al., ``Triton-distributed: Programming overlapping kernels on
distributed AI systems with the Triton compiler,'' arXiv:2504.19442, 2025.

\bibitem{mltriton}
D. Wang, W. Zhu, L. Ling, E. Tiotto, Q. Wang, W. Tsang, J. Opperman, and
J. Deng, ``ML-Triton, a multi-level compilation and language extension to
Triton GPU programming,'' arXiv:2503.14985, 2025.
\end{thebibliography}

\end{document}
